\documentclass[12pt,a4paper]{article}

\usepackage[a4paper,margin=1in]{geometry}
\usepackage{hyperref}

\title{Computer Science Portfolio Website Report}
\author{ Antony Joshi  (GH1056921)}
\date{\today}

\begin{document}

\maketitle
\tableofcontents
\newpage

\section{Introduction}

The implementation of the report documents shows the design and content of my computer science portfolio website developed as part of the B201 computer science lab module. The portfolio shows the professional showcase of my technical skills, education and projects for potential employers reviewing my application for a working student position.

These days having a properly designed online presence is really important for a computer science student in job market and portfolio websites provides potential employers can learn about a candidate's background, skills, and achievements.

There is a main objective for this project is to create portfolio that shows my qualification and practical competencies in web development. The portfolio website is hosted on GitHub Pages and integrates various technologies learned throughout the module.


Portfolio Website:

\url{https://antonyjoshi2006.github.io/CS-Assignment-/}

Git Hub repository:

\url{https://github.com/antonyjoshi2006/CS-Assignment-}

This report will explain the structure of the portfolio website and its corresponding Github repository.

\section{Portfolio Website Structure}

There is a six main section for portfolio website and each serving a specific purpose which are as follows

Main section: As the landing area the Main section makes a first impression on viewers so it includes my bio like my name and a quick description of myself as a tech enthusiast and student of computer science. There are two conspicuous call-to-action buttons that direct people to view my work and contact me.

About Section: In about section we can see my personal introduction that explains my background, interest in technology and my professional aspirations. This section humanizes the portfolio by sharing my motivations and goals as well as helping employers to understand not just my skills but also my personality and my work ethics.

Education Section: This section states my B.Sc. Computer Science program at GISMA University of Applied Sciences. This shows my areas of interest which are as follows- programming, software development emerging technologies and provides my expected graduation date of 2028.


Skill Section: The Skills section presents my abilities in three categories: Technical skills such as python, HTML, SQL, Web Development and soft skills are problem-solving,continuous learning collaboration as well as interests like developing technologies,real world solutions.This approach allows viewers to easily assess my qualifications from multiple ways.

Portfolio Section: The portfolio section is designed to see my completed projects with descriptions, screenshots and links to the source code.Currently this section indicates that projects are coming soon as I am still e in my academic journey.

Contact Section: In this section the email address was provided so it will helps to invite professional connections. This is done so that I can express my openness to opportunities,collaborations and networking.

 Page Layout

The website is made with clean and modern layout with proper styling.Every section occupies its own visual block which allows users to scroll through content naturally and there is a fixed direction menu that helps the viewers go to specific sections.

\section{GitHub Repository Structure}

The Git hub have clear sepration between different components of the project. A repository is important not only for personal but also for potential employers often review candidates, Github profile helps to assess the coding practices


\subsection{Directory Structure}

Directory Structure

The repository contains the following folders and files:

ASSETS: Contains all image and media files used on the website 

CV: Contains the Latex source file and compiled PDF of my curriculum vitae 

REPORT: Contains this project report in Latex format 

Index.html: The main portfolio webpage 

README.md: Repository documentation explaining the project 

LICENSE: The license file for the project

\end{itemize}

\subsection{Folder Description}

In assets folder integrateds all media files in the website including decorative elements. Keeping all media in a single folder simplifies maintenance because file paths remain consistent and updating.

In cv folder contains my curriculum vitae created using Latex. The CV follows a professional template and includes my contact information, skills, education, experience,activities,awards and languages. Having the CV in both Latex source format and compiled PDF demonstrate proficiency in professional document preparation.

The REPORT folder contains this project report, also prepared using LaTeX. Keeping the report separate from other materials ensures the repository remains organized.


\subsection{Documentation and Version Control}

The repository's documentation is provided by the README.md file, which explains the project goal, the folder structure and how to access the live websiter epository may understand how the project progresses thanks to the repository's meaningful commit history and comprehensive messages.

\section{DESIGN DECISIONS}

component of the portfolio was created with particular goals and explanations in mind.

\subsection{Single page design}

I chose single page design over than a multipage website for several important reason:

\begin{itemize}
    \item Simplicity: A single html file is easier to see and maintain
    \item User Experience: visitors can view all content without looking different pages which decreases friction and the likelihood that they will depart before seeing crucial information.
    \item Mobile Compatibility:Single page designs perform better on mobile device.
\end{itemize}

\subsection{Visual Design Choices}

The color scheme is used professional and create a clean aesthetic for the portfolio. The design choices were guided by several principle:

To create a distinct hierarchy, typography was carefully create a distinct hierarchy, typography was carefully examined. While body text utilize a legible size and weight headings employ larger, bolder fonts to stand out. Visitors can quickly scan the page and find portions of interest thanks to this hierarchy.

White space was help throughout the design. Suffcient spacing between elements improves readability and prevents the page from feeling cluttered.Padding is used in each part to visually distinguish it from neighboring ones.Consistency was maintained across all elements such as buttons, links and navigation items followthe same patterns and creating a cohesiv experience.


\subsection{Content Organisation Strategy}

The about section provides context and humanizes the portfolio 

The education section demonstrates qualifications 

The skills section highlights specific capabilities 

The portfolio section offers practical demonstrations 

The contact section enables followup and connection

This feels guides visitors through a logical narrative from introduction to action.

\section{Tools and Technologies}

This section clarifies the tools and technologies is used to develop the portfolio website.


 Web Development Technologies

The materials of the portfolio was organized using HTML also enhanced by HTML's header, navigation, section and footer components. 

CSS was used for layout and responsive design .CSS shows the visual presentation of the website, including colors,spacing and positioning.

Github Hosting option was chosen because Github Pages offers dependable and free hosting straight from the repository and promptly deploys updates as changes are pushed and this gives a professional URL format and doesn't require server configuration.

\subsection{Document preparation}

Latex was used for creating both CV and this project report.LaTeX is a typesetting system widely used in academic and scientific publishing.It provides many advantages such as professional typography, consistent formatting through documents etc...


\subsection{Version Control}

GIT was used for tracking all development changes.Git enables developers to track changes and revert to prior versions if needed.Every update to the portfolio was ensured with descriptive notes creating a complete history of the development process.

Git hub was used for remote repository,collaboration tools also integration with GitHub Pages for automatic deployment and visibility for employers who might look over my profile 
\section{Reflection}

This part offers a critical assessment of the portfolio pointing out its advantages and flaws as well as potential improvement strategies.

The website is easy to navigate and has a clear structure.Using the navigation menu or by scrolling viewers can locate information quickly.The single-page layout maintains a balance between thoroughness and simplicity.

The design exudes professionalism that is suitable for applications for jobs in the technology industry. A sense of skill and attention to detail is created by the layout,font and color scheme.

The repository is properly sorted with commit changes and README.md documentation. This demonstrates professional development practies that employers look for.

Using LaTeX for the CV and this report, shows technical writing skills particularly in computer science and academic contexts.

\subsection{Flaws}

The portfolio section currently lacks completed projects to showcase. As a student early in academic journey I am still working on my projects and once they are completed they will be updated.This is the most significant weakness in portfolios.

The website uses only HTML and CSS without JavaScript for improved functionality.

Contact is limited to an email link which requires viewers to open their email client.


The website uses only HTML and CSS without JavaScript for improved functionality.

Contact is limited to an email link which requires viewers to open their email client.

\section{Conclusion}

This portfolio project has provided practical experience in web development.The process of planning as well as designing and implementing a complete website reinforces concepts learned throughout the module and demonstrates how these skills apply to real-world scenarios.The resulting portfolio website works as a foundation for showcasing my technical skills to the employers.While it currently reflects my status as an early stage student the structure accommodates the growth as I complete more projects .

The development process reinforces several important lessons thoughtful design choices profit better results than random ones,proper documentation is an essential professional practice and each and every project should contain continuous improvement from the outset.These skills are instantly applicable to professional software development practice which will help me in my computer science studies and jobs.

\section{Refrences}

Github: \url{https://github.com/Manciukic/MAltaCV}

Riccardo Mancini

\end{document}
```
